% Source: https://artofproblemsolving.com/community/c3h2835479_is_this_a_valid_alternate_solution
% Stable identity: aops-topic-2835479-q01
Let the incircle touch $BC$, $CA$ at $X$, $Y$. Tangent segments from a point are equal: $AZ=AY=x$, $BZ=BX=y$, $CX=CY=z$.
\begin{asy}
usepackage("amsmath");
size(4.5cm);
defaultpen(fontsize(10pt));
pair A=(0,0),B=(6,0),C=(1.5,4);real a=abs(B-C),bb=abs(C-A),c=abs(A-B);pair I=(a*A+bb*B+c*C)/(a+bb+c);
pair Z=(I.x,0);pair Y=A+(abs(Z-A))*unit(C-A);pair X=B+(abs(Z-B))*unit(C-B);
draw(A--B--C--cycle);draw(circle(I,I.y));dot(Z);dot(Y);dot(X);
label("$A$",A,SW);label("$B$",B,SE);label("$C$",C,N);label("$Z$",Z,S);label("$Y$",Y,W);label("$X$",X,E);
\end{asy}

Then $x+y=c$, $y+z=a$, $z+x=b$, so $AZ=x=\boxed{\dfrac{b+c-a}2}$.

The proposed formula fails, e.g. for an equilateral triangle of side $s$: the true value is $\dfrac s2$, but the formula gives $\dfrac{60}{90}\cdot\dfrac s2=\dfrac s3$.
